\subsection{Questão 58}

\subsubsection{Enunciado}

Use a aproximação do Problema 55 para calcular

\begin{enumerate}
    \item[(a)] a altitude em relação à superfície na qual o módulo do campo magnético da Terra é \(50{,}0\%\) do valor na superfície na mesma latitude;

    \item[(b)] o módulo máximo do campo magnético na interface do núcleo com o manto, \(2900\,\text{km}\) abaixo da superfície da Terra;

    \item[(c)] o módulo e

    \item[(d)] a inclinação do campo magnético na Terra no pólo norte geográfico.

    \item[(e)] Explique por que os valores calculados nos itens (c) e (d) não são necessariamente iguais aos valores medidos.
\end{enumerate}



\vspace{1cm}
